% ===========================================================================
%  CTD validation of novel predictions -- Results section
%  Requires: booktabs, graphicx, hyperref (load hyperref LAST in the preamble).
%  \TODO{} items to resolve before submission.
% ===========================================================================
\providecommand{\TODO}[1]{\textbf{[TODO: #1]}}

% ------------------------------------------------------------------ METHODS
\subsection{Validation against curated therapeutic records}
\label{sec:ctd-methods}

We validated the highest-scoring unobserved pairs against the Comparative
Toxicogenomics Database (CTD; release \TODO{YYYY-MM-DD}), keeping only curated
relationships annotated as \emph{therapeutic}. Because CTD indexes by MeSH while
Cdataset uses DrugBank accessions and OMIM MIM numbers, both axes were mapped
($657/658$ drugs; \TODO{N}$/409$ diseases; Supplementary Note~\TODO{S}). Each
method was fitted to the complete matrix with its optimal hyper-parameters, and
all entries that were zero in the training matrix were ranked by predicted
probability. Following Zhang \emph{et al.}~\cite{drimc} we evaluated the ranking
at two levels: \emph{ranked by pairs}, the fraction of the top $N$ pooled pairs
supported by a curated therapeutic record absent from the training matrix; and
\emph{ranked per drug}, the fraction of drugs with at least one such record among
their own top-$K$ diseases. As in DRIMC, the per-drug denominator is restricted
to the $361$ of $658$ drugs ($54.9\%$) for which such a record exists. All
methods used an identical pipeline, so the comparison isolates the ranking each
model produces.

% ------------------------------------------------------------------ RESULTS
\subsection{Validation of novel drug--disease associations}
\label{sec:ctd-results}

Ranked by pairs (Table~\ref{tab:ctd_pairs}, Figure~\ref{fig:ctd}a), BVSIMC
verified $11.0\%$ of its top $500$ pooled predictions, against $8.4\%$ for SGIMC,
$5.6\%$ for DRIMC and $4.6\%$ for IMC, and led at every depth to $N=1500$
($9.1\%$ versus $7.4\%$)---a relative improvement of $23$--$37\%$. Precision fell
with $N$ for BVSIMC and SGIMC but rose for IMC and DRIMC, so only the former two
concentrate verified associations at the very top of the ranking, which is what
matters when a short candidate list is carried forward for experimental
follow-up.

Ranked per drug (Table~\ref{tab:ctd_drugs}, Figure~\ref{fig:ctd}b), BVSIMC was
highest at every cut-off, reaching $34.3\%$ at top-5 and $44.3\%$ at top-10
against $29.1\%$ and $39.3\%$ for the closest competitor; no single method was
consistently second. The margin widens with $K$, from $1.1$ percentage points at
top-1 (about four drugs) to $5.2$ and $5.0$ points at top-5 and top-10 (about
$19$ and $18$ drugs). \TODO{McNemar's exact test on the paired per-drug outcomes;
at top-1 and top-3 the gap is only 4--5 drugs, so write \emph{comparable to}
there unless the test says otherwise} Absolute rates are bounded by the
benchmark: only $54.9\%$ of Cdataset drugs have any curated therapeutic record
outside the training matrix, against $83.2\%$ of LRSSL drugs in the original
DRIMC evaluation, as Cdataset's OMIM diseases are predominantly rare Mendelian
phenotypes. For reference, a degree-product baseline using no side information
recovered $0.0$--$1.9\%$ across the same cut-offs.

Table~\ref{tab:ctd_cases} illustrates the recovered indications for the six
drugs with the most CTD-verified predictions. Of the \TODO{30} pairs shown,
\TODO{n} carry a curated therapeutic record absent from the training matrix,
including \TODO{drug A} for \TODO{disease} and \TODO{drug B} for \TODO{disease}.
These drugs were selected for illustration rather than at random, so the table
shows what a successful prediction list looks like and not the average case; the
rates in Tables~\ref{tab:ctd_pairs} and~\ref{tab:ctd_drugs} are the
representative figures.

Taken together, BVSIMC ranked verified associations highest of the five methods
under both criteria, and by the larger and more stable margin at the pair level.
Its precision also decays with depth rather than rising, so the pairs it places
at the very top are the ones most likely to be real---the property that decides
whether a shortlist is worth pursuing experimentally. Absolute rates remain low
in absolute terms, but this reflects the benchmark rather than the models: with
$54.9\%$ of Cdataset drugs having any verifiable indication at all, and CTD
holding few therapeutic records for the rare Mendelian phenotypes that dominate
its disease axis, the achievable ceiling on this dataset is well below that of
LRSSL.

% ------------------------------------------------------------------- TABLES
\begin{table}[htbp]
\centering
\caption{\textbf{Ranked by pairs.} Percentage of the top $N$ drug--disease pairs,
pooled across all drugs, supported by a CTD-verified therapeutic record
(Cdataset). Best in bold.}
\label{tab:ctd_pairs}
\begin{tabular}{lrrrr}
\toprule
Method & $N=500$ & $N=800$ & $N=1000$ & $N=1500$ \\
\midrule
BVSIMC & \textbf{11.0} & \textbf{10.8} & \textbf{10.0} & \textbf{9.1} \\
SGIMC  &  8.4 &  8.0 &  7.3 &  7.4 \\
IMC    &  4.6 &  5.4 &  5.8 &  5.7 \\
DRIMC  &  5.6 &  5.8 &  6.5 &  6.8 \\
NRLMF  &  1.4 &  1.2 &  1.4 &  1.3 \\
\bottomrule
\end{tabular}
\end{table}

\begin{table}[htbp]
\centering
\caption{\textbf{Ranked per drug.} Percentage of eligible drugs with at least one
CTD-verified therapeutic indication, absent from the training matrix, among their
top-$K$ predicted diseases (Cdataset, $n=361$). Best in bold.}
\label{tab:ctd_drugs}
\begin{tabular}{lrrrr}
\toprule
Method & top-1 & top-3 & top-5 & top-10 \\
\midrule
BVSIMC & \textbf{13.0} & \textbf{25.5} & \textbf{34.3} & \textbf{44.3} \\
SGIMC  & 11.9 & 23.5 & 28.8 & 38.2 \\
IMC    & 10.0 & 24.1 & 29.1 & 38.2 \\
DRIMC  &  6.9 & 18.3 & 27.1 & 39.3 \\
NRLMF  &  1.4 &  3.6 &  4.4 &  7.8 \\
\bottomrule
\end{tabular}
\end{table}

% ------------------------------------------------------------------- FIGURE
\begin{figure}[htbp]
\centering
\includegraphics[width=\textwidth]{figures/cdataset_methods.pdf}
\caption{Validation of novel drug--disease predictions against CTD curated
therapeutic records (Cdataset). (\textbf{a})~\textbf{Ranked by pairs}: percentage
of the top $N$ pooled pairs supported by a CTD-verified therapeutic record.
(\textbf{b})~\textbf{Ranked per drug}: percentage of drugs with at least one such
record among their top-$K$ predicted diseases ($n=361$ eligible drugs). All
methods were fitted to the complete matrix with their optimal hyper-parameters
and evaluated through an identical pipeline. CTD release \TODO{YYYY-MM-DD}.}
\label{fig:ctd}
\end{figure}

% -------------------------------------------------------- SUPPLEMENTARY NOTE
\subsection*{Supplementary note: identifier mapping}

Drug mapping resolved $657/658$ DrugBank accessions: $491$ by CAS registry
number, $135$ by preferred name, $9$ by synonym and $22$ after removing salt or
ester qualifiers, which are needed because DrugBank names the parent compound
while CTD indexes the marketed salt. Disease mapping resolved \TODO{N}$/409$ OMIM
entries: $295$ through CTD's curated OMIM cross-references and $97$ through exact
matches of a full OMIM title---preferred or alternative---against a CTD disease
name or synonym; five matched only after removing an OMIM qualifier, of which two
were retained and three discarded as over-broad. Unresolved entries were kept in
the rankings but treated as unverifiable, which can only lower a reported rate.
Because CTD does not resolve genetic subtypes, the mapped entries collapse onto
\TODO{378} distinct concepts, but duplicate (drug, concept) entries account for
under $1\%$ of per-drug top-10 slots and $2\%$ of the pooled ranking, so no
de-duplication was applied.

% -------------------------------------------------- DATA & CODE AVAILABILITY
\subsection*{Data and code availability}

Source code and datasets are available at
\href{https://github.com/Sijianf/BVSIMC}{\nolinkurl{github.com/Sijianf/BVSIMC}}.

% ===========================================================================
%  BEFORE SUBMISSION
%  1. Run McNemar's exact test; top-1/top-3 gaps are only 4-5 drugs of 361.
%  2. State which disease-mapping criterion was used; show the other two in SI.
%  3. NRLMF sits well below the others -- report K1/K2 (published defaults were
%     tuned on 54-445 drug datasets, not 658) so readers can judge.
%  4. Consider a case-study table (cf. DRIMC Table 4), marking CTD-verified and
%     literature-supported pairs separately.
% ===========================================================================
